% trajectories.tex -- Methodology: payload trajectories and computation (draft 4 Oct 2026).
% Definitions: mpc_planner/load_trajectory.py. Solve times measured from the planner logs
% (solve_ms column): rig 1 Oct flights h5_z100_rot120, c1_z100, f8_z100 (three UAVs, laptop,
% no Gazebo); Gazebo R0978/R0979 (three UAVs), R0931/R0935 (four UAVs), laptop at RTF ~0.3.

\subsection{Payload trajectories}
\label{sec:trajectories}

The payload reference is the lift schedule in height plus a horizontal offset from a trajectory.
The offset starts once the lift has reached the target height, and is applied to the whole formation,
so the UAVs keep their places around the payload. The planner evaluates the trajectory at the time
of every horizon node, so the plan sees the next two seconds of the path. The acceleration
$\mathbf{a}_L^{\mathrm{ref}}$ used for the cable references (Section~\ref{sec:controller}) is
its analytic second derivative.

Trajectories that start and end at rest use a smootherstep sweep of the path angle,
\begin{equation}
    \theta(t) = 2\pi\, s\!\left(\tfrac{t}{T}\right), \qquad s(u) = 6u^5 - 15u^4 + 10u^3, \quad u \in [0,1],
\end{equation}
whose first and second derivatives vanish at both ends, so the manoeuvre has no step in velocity or
acceleration that the suspended payload would take as a kick. The duration $T$ is chosen so that the
peak tangential speed equals the commanded speed $v$. The kinds are:
\begin{description}
    \item[Hover.] No horizontal motion; the payload holds the target height until LAND.
    \item[Line.] A shuttle along $x$, $\Delta x = \tfrac{d}{2}\left(1 - \cos\tfrac{2vt}{d}\right)$, repeated until LAND.
    \item[Circle.] One lap, $\Delta x = r\sin\theta$, $\Delta y = r(1-\cos\theta)$, with
    $T = 1.875 \cdot 2\pi r / v$. It starts at the lift point heading $+x$ and ends there at rest.
    \item[Figure-eight.] One Gerono lemniscate, $\Delta x = a\sin\theta$, $\Delta y = \tfrac{a}{2}\sin 2\theta$,
    with $T = 1.875 \cdot 2\pi a\sqrt{2} / v$ (the peak speed is at the crossing).
    \item[Spin.] The circle with one full turn of the payload's yaw, synchronised with the path.
    \item[Orbit.] The circle flown continuously at constant speed until LAND. The angular rate
    $v/r$ is reached through a smootherstep ramp over 4\,s, and the angle is its exact integral,
    so position, velocity and acceleration are continuous. The orbit keeps the payload moving
    for as long as a reconfiguration takes.
\end{description}

The results use an orbit of radius 0.5\,m at 0.125\,m/s (one lap in 25\,s) and a figure-eight of
$a = 1.0$\,m at 0.25\,m/s (67\,s, a $2 \times 1$\,m path), sized for the 6\,$\times$\,6\,m flight area
with the UAVs about 0.7\,m outside the payload's path.

\paragraph{Computation.} Both levels run on the same laptop as the motion-capture and radio nodes.
Table~\ref{tab:solve_times} gives the planner's solve time per tick. On the rig, with three UAVs,
the median is 24--30\,ms and the 95th percentile 49--55\,ms, inside the 60\,ms budget and the
100\,ms tick. In Gazebo the same laptop also runs the physics and the simulated sensors, which slows
the solves, but the simulation runs at about 0.3 of real time, so the planner still meets its tick in
simulated time. The four-UAV planner has not yet been timed on the rig.

\begin{table}[h]
\centering
\caption{Planner solve time per 10\,Hz tick (wall clock, one laptop).}
\label{tab:solve_times}
\begin{tabular}{llccc}
\toprule
Platform & Case & Median (ms) & 95th pct. (ms) & Max (ms) \\
\midrule
Rig, 3 UAVs & hover at 1.0\,m & 30.4 & 54.9 & 73.3 \\
            & orbit & 24.6 & 51.9 & 71.8 \\
            & figure-eight & 23.8 & 49.3 & 76.6 \\
\midrule
Gazebo, 3 UAVs & hover at 1.0\,m & 46.7 & 97.8 & 137.9 \\
               & orbit & 47.1 & 91.0 & 149.6 \\
Gazebo, 4 UAVs & hover at 1.0\,m & 118.5 & 193.2 & 298.5 \\
               & figure-eight & 106.2 & 189.1 & 313.7 \\
\bottomrule
\end{tabular}
\end{table}
